% STATUS: VERIFIED / AUDITED
% AUDIT FILE: audits/A11_hopf_spectral_conditions.tex
% AUDIT VERDICT: VERIFIED
% LAYER VERDICTS: general frequency/strict candidate VERIFIED;
%   simplicity/eigenvector/resolvent VERIFIED; candidate surface VERIFIED;
%   realizability criterion VERIFIED; five-node candidate VERIFIED;
%   scope/P12-P13 handoff VERIFIED; OVERALL VERIFIED.
% REQUIRED CORRECTIONS: NONE; no mathematical or wording correction.
% AUDITED VERSION: P11 v1.0; 2026-10-03; 25914 bytes; 600 lines.
% AUDITED FULL-FILE SHA-256:
%   4be8e9b1e64c022b635f3cca9f7f49ffbc5358a2250dda1715cb60e82d9b073b
% AUDITED MATHEMATICAL-BODY SHA-256:
%   faa2341cc2dce31b00f093aca4c33566d02b261dbd8211ed602c0ce29f837fce
% CLOSURE DATE: 2026-10-04
% CURRENT VERSION: P11 v1.1; metadata-only closure; no mathematical revision.
% CURRENT FULL-FILE HASH: recorded separately in project_manifest.md and
%   audit_closure_report.txt to avoid a self-referential full-file hash.
% CURRENT MATHEMATICAL-BODY SHA-256:
%   faa2341cc2dce31b00f093aca4c33566d02b261dbd8211ed602c0ce29f837fce
% BODY HASH SCOPE: bytes from first line-anchored section command through EOF.
% PRESERVATION: mathematical body byte-identical to exact audited v1.0.
% STATUS NOTE: historical DRAFT/audit-pending body prose is preserved;
%   this header and the current manifest govern lifecycle status.
% VALIDATION PROVENANCE: A11 reconstructed the mathematics analytically
%   and supplied independent exact/finite corroboration. Developer validation
%   was authenticated/inventoried, not rerun or used as proof.
% PLANNING: clarification accepted in open_decisions.md; original proposal
%   and all frozen documents unchanged; recommendation separate from verdict.
% DOWNSTREAM: P12 CLEARED TO BEGIN, but not begun by this closure.
% DEPENDS ON: P02, P04, P05 VERIFIED / AUDITED
% CONDITIONAL SPECIALIZATION INPUT: P07 VERIFIED / AUDITED
% MODEL: canonical DM, q=2
% ROLE: Hopf spectral conditions and spectral simplicity
% FIVE-NODE CANDIDATE: n=5, q=2
% NUMERICAL DEPENDENCIES: none
% NEXT DEPENDENCY: P12
% LATER DEPENDENCIES: P13, P14
% EXCLUSIONS: no transversality, no Lyapunov coefficient, no Hopf theorem,
%             no periodic branch, no rotating-wave theorem
% CONTRACT: 00_canonical_model.tex; 01_notation_and_assumptions.tex;
%           shared_macros.tex. P03 inherited through P04; no new DAG edge.
% CLAIMS: C072 and spectral part of C074; P11-D1--P11-D8.
% C073: archived numerical candidate remains unreproduced and unpromoted.
% BASELINE: complete authenticated cyclic_network_project_P10_audited.zip.

\section{P11: Hopf spectral conditions and spectral simplicity}
\label{sec:P11}
\noindent\textbf{Status: DRAFT; independent A11 audit pending.}
We prove spectral assertions for the canonical demand-modulated model at
positive synchrony. Direct inputs are P02, P04 and P05; P07 is used only
for the five-node specialization. P03's coefficients are inherited through
P04. P08--P10 and P12 or later results supply no mathematical premise.
In particular, an imaginary pair here is only a \emph{spectral Hopf
candidate}. No parameter crossing, nonlinear nondegeneracy, periodic
orbit or propagation theorem is asserted. The physical path is frozen for
P12 without differentiating along it. The optional realizability result
below concerns spectral boundaries, not first instability.

\subsection{Setting, representative classes and frequency}
\label{slot:P11-1}
Fix $q=2$, $n\ge5$, $a,k,c,g,h>0$, and a specified positive synchronous
equilibrium $x^*\one$. Set
\begin{equation}
 s=1+x^*,\qquad C=\frac cs>0,\qquad
 G=\frac{gx^*}{s^2}>0,\qquad H=hx^*>0.
 \label{eq:P11-effective}
\end{equation}
Use P04's verified spectrum and unitary diagonalization:
\begin{align}
 \lambda_m&=d+Ce^{-\ii\theta_m}+Ge^{2\ii\theta_m}
                       -He^{-2\ii\theta_m},\quad \theta_m=2\pi m/n,
                       \label{eq:P11-input}\\
 r_m&=d+C\cos\theta_m+(G-H)\cos2\theta_m,\nonumber\\
 \omega_m&=-C\sin\theta_m+(G+H)\sin2\theta_m,\nonumber\\
 J&=\Fourier\diag(\lambda_m)\Fourier^*,\qquad
 \lambda_{n-m}=\overline{\lambda_m}\quad(1\le m<n).\nonumber
\end{align}
Here $J=DF^{\DM}(x^*\one)$; $\omega_m$ is a signed imaginary part.
The reserved positive frequency is chosen separately when needed.

Let $N=\lfloor n/2\rfloor$, $u_j=\cos(2\pi j/n)$ and
$L_j$ be P05's representative label classes: $L_0=\{0\}$,
$L_N=\{n/2\}$ if $n$ is even, and otherwise
$L_j=\{j,n-j\}$. Write
\[
 \mathcal S=\argmax_{0\le j\le N}r_j
\]
for the \emph{complete} selected representative set, including every tie.
The selected Fourier labels are $\bigcup_{j\in\mathcal S}L_j$.
This notation distinguishes $\mathcal S$ from the scalar $G+H$.

\begin{lemma}[Reachable nonuniform representatives are genuine pairs]
\label{lem:P11-pairs}
Every reachable nonuniform representative satisfies
$1\le j\le j_{\max}:=\lfloor(n+1)/4\rfloor<n/2$.
Consequently $0<\theta_j<\pi$, $\sin\theta_j>0$, and
$L_j=\{j,n-j\}$ has two different labels. No even-$n$ Nyquist
singleton is included.
\end{lemma}
\begin{proof}
P05 Theorem~\ref{thm:P05-reachable} supplies the reachable list.
Since $(n+1)/4<n/2$ for $n>1$, its nonzero members satisfy the
claimed strict bounds. The angle and sine statements follow.
This counts Fourier labels only, not numerical eigenvalue multiplicities.
\end{proof}

\begin{proposition}[Exact signed frequency and its zero locus]
\label{prop:P11-frequency}
For any genuine nonuniform pair representative $0<j<n/2$,
\begin{equation}
 \boxed{\omega_j=\sin\theta_j\,[2(G+H)u_j-C].}
 \label{eq:P11-frequency}
\end{equation}
Thus $\omega_j=0$ if and only if $2(G+H)u_j=C$.
If $u_j\le0$, then $\omega_j<0$ automatically. If $u_j>0$,
the zero-frequency locus is exactly
\begin{equation}
 G+H=\frac{C}{2u_j}.
 \label{eq:P11-zero-frequency}
\end{equation}
For positive $u_j$ the sign of $\omega_j$ is the sign of
$G+H-C/(2u_j)$. These are spectral statements, not bifurcation classifications.
\end{proposition}
\begin{proof}
Insert $\sin2\theta_j=2\sin\theta_j\cos\theta_j$ in
\eqref{eq:P11-input}. The factor $\sin\theta_j$ is strictly positive.
If $u_j\le0$, the bracket is at most $-C<0$; otherwise divide by
$2u_j>0$. The Nyquist singleton is outside the proposition's scope.
\end{proof}

\subsection{Strict candidates, simplicity and the stable complement}
\label{slot:P11-2}
\begin{definition}[Strict nonuniform spectral Hopf candidate]
\label{def:P11-candidate}
A positive synchronous equilibrium of the canonical $q=2$ model is a
\emph{strict nonuniform spectral Hopf candidate in representative $j$}
if $j\ne0$, $\mathcal S=\{j\}$, $r_j=0$, and $\omega_j\ne0$.
Uniqueness refers to representatives and includes the exclusion of all
maximal ties. This definition does not assert a Hopf bifurcation.
\end{definition}

\begin{theorem}[General strict spectral criterion]
\label{thm:P11-strict}
At a strict candidate, the entire imaginary-axis spectrum is exactly
\begin{equation}
 \lambda_j=\ii\omega_j,\qquad
 \lambda_{n-j}=-\ii\omega_j,\qquad \omega_j\ne0.
 \label{eq:P11-critical-pair}
\end{equation}
Each value is algebraically and geometrically simple over $\C$ and the
two values are distinct. Every other Fourier label has $r_\ell<0$.
The number
\begin{equation}
 \eta=-\max_{\ell\notin L_j}r_\ell
 \label{eq:P11-gap}
\end{equation}
is strictly positive, and all noncritical eigenvalues satisfy
$\Ree\lambda_\ell\le-\eta$. The real critical space has dimension
two and its real invariant complement has dimension $n-2$.
Moreover
\begin{equation}
 0,\ 2\ii\omega_j\notin\spec(J),\qquad
 J^{-1}\ \hbox{and}\ (2\ii\omega_j I-J)^{-1}\ \hbox{exist}.
 \label{eq:P11-resolvent-gates}
\end{equation}
\end{theorem}
\begin{proof}
P05 reachability and Lemma~\ref{lem:P11-pairs} give two distinct
labels. P04 conjugacy proves \eqref{eq:P11-critical-pair}.
For each other representative $t$, complete strict selection gives
$r_t<r_j=0$: equality would place $t$ in $\mathcal S$.
Every remaining Fourier label belongs to one of those other classes,
so every remaining real part is negative, including the uniform label
and any Nyquist label.

If another label outside $L_j$ produced $\lambda_j$, it would have
real part zero, contradicting that strict inequality. The other label
$n-j$ cannot produce $\lambda_j$ either, since $\omega_j\ne0$.
Thus precisely one label produces $\lambda_j$, and the same argument
applies to $\lambda_{n-j}$. P04
Proposition~\ref{prop:P04-multiplicity} equates both complex
multiplicities with the number of coincident labels. This proves
simplicity, rather than inferring it merely from the Fourier basis.
The real plane and complement follow from the conjugate Fourier
spaces; an explicit orthonormal real basis is given below.

There are $n-2\ge3$ noncritical labels, a nonempty finite set with
strictly negative real parts. Its maximum is negative, proving
$\eta>0$ and the claimed bound. Neither zero nor $2\ii\omega_j$
can equal a noncritical eigenvalue because their real parts are zero.
Neither equals $\pm\ii\omega_j$: the first exclusion uses
$\omega_j\ne0$, and the second would require $2=1$ or $2=-1$
after division by $\ii\omega_j$. This proves both inverse gates.
\end{proof}

\begin{proposition}[Exact real and adjoint normalization]
\label{prop:P11-vectors}
At the strict candidate, choose
\begin{equation}
 v_H=v^{(j)},\qquad p_H=v_H,\qquad
 e_c=\sqrt2\,\Ree v_H,\quad e_s=\sqrt2\,\Imm v_H.
 \label{eq:P11-vectors}
\end{equation}
Under the project's Hermitian pairing $\inner{p}{v}=p^*v$,
\begin{equation}
 Jv_H=\ii\omega_jv_H,\qquad
 J^*p_H=-\ii\omega_jp_H,\qquad \inner{p_H}{v_H}=1.
 \label{eq:P11-adjoint}
\end{equation}
The real vectors $e_c,e_s$ are orthonormal and span the real critical
plane. In the ordered basis $(e_c,e_s)$ the matrix is
\begin{equation}
 \left.J\right|_{\rm crit}=
 \begin{pmatrix}0&\omega_j\\-\omega_j&0\end{pmatrix}.
 \label{eq:P11-real-block}
\end{equation}
In its coefficient plane the signed angular velocity is $-\omega_j$.
\end{proposition}
\begin{proof}
P04 gives $J=\Fourier\Lambda\Fourier^*$, hence
$J^*=\Fourier\overline\Lambda\Fourier^*$. The same normalized
Fourier column therefore solves the opposite-sign adjoint equation.
This agrees with Section~\ref{sec:notation}, including conjugate
linearity in the first argument.

Write $v_H=A+\ii B$ with real vectors $A,B$.
The columns $v_H$ and $\overline{v_H}=v^{(n-j)}$ are orthogonal
and each has norm one. Thus $v_H^{\mathsf T}v_H=0$ and
$v_H^*v_H=1$, giving
$\|A\|^2=\|B\|^2=1/2$ and $A^{\mathsf T}B=0$.
Every real vector in the complex critical span is a real linear
combination of $A,B$, proving the plane assertion. Taking real and
imaginary parts of $J(A+\ii B)=\ii\omega_j(A+\ii B)$ gives
$Je_c=-\omega_j e_s$ and $Je_s=\omega_j e_c$, which are the
columns of \eqref{eq:P11-real-block}. For coefficients $(z_c,z_s)$,
$\dot z_c=\omega_jz_s$ and $\dot z_s=-\omega_jz_c$.
This is a linear orientation statement only. If $\omega_j<0$, one
may instead choose label $n-j$ to obtain the positive frequency
$|\omega_j|$; here the signed representative convention is retained.
\end{proof}

\begin{remark}[Ties and zero frequency are excluded separately]
\label{rem:P11-degeneracies}
If two different representative classes have $r_j=r_t=0$, their
labels give additional imaginary-axis values or larger multiplicities.
Either possibility fails the single-simple-pair-with-stable-complement
hypothesis. A tie of real parts alone is not equality of numerical
eigenvalues. In particular choosing one member of a maximal tie does
not satisfy Definition~\ref{def:P11-candidate}.

If a genuine pair has $r_j=\omega_j=0$, both labels produce the
same real zero eigenvalue. P04 makes it semisimple, with complex
algebraic and geometric multiplicity exactly
$|\{m:\lambda_m=0\}|\ge2$. Under strict selection of that
representative its multiplicity is exactly two and its real kernel
has dimension two. Without strict selection other zero labels must
also be counted. This is ineligible for the simple Hopf package.
No named multiple-criticality or steady-state bifurcation is classified.
The two-dimensional real plane in Theorem~\ref{thm:P11-strict}, by
contrast, corresponds to two distinct simple complex eigenvalues.
\end{remark}

\subsection{Regular spectral surfaces and physical admissibility}
Fix a reachable nonuniform representative $1\le j\le j_{\max}$
throughout this subsection, so $0<\theta_j<\pi$.
Use P02's nonzero equilibrium equation in branch coordinates:
\begin{equation}
 a=(k+h)x-\frac{c-g}{1+x}>0,\qquad x,k,c,g,h>0.
 \label{eq:P11-chart}
\end{equation}
This is exactly P05's open chart domain $\mathcal U_{\rm eq}$.
Combining P04's diagonal coefficient with the equilibrium equation gives
\[
 d=a-2kx-\frac{cx}{(1+x)^2}-\frac g{1+x}-hx
   =-kx-\frac{c(1+2x)}{(1+x)^2}.
\]
For fixed representative $j$, put $v_j=\cos2\theta_j=2u_j^2-1$;
this scalar $v_j$ is not a Fourier vector $v^{(j)}$. Define
\begin{equation}
 \begin{split}
 R_j(x,k,c,g,h)={}&-kx-\frac{c(1+2x)}{(1+x)^2}
                 +\frac{c}{1+x}u_j\\
                &+x\left(\frac{g}{(1+x)^2}-h\right)v_j.
 \end{split}
 \label{eq:P11-R}
\end{equation}
Substitution in P04's real dispersion shows exactly $R_j=r_j$.

\begin{proposition}[Regular spectral boundary and excluded frequency locus]
\label{prop:P11-surface}
At fixed $x,c,g,h$,
\begin{equation}
 \frac{\partial R_j}{\partial k}=-x<0,
 \qquad
 k_H=\frac1x\left[-\frac{c(1+2x)}{(1+x)^2}
       +\frac{c}{1+x}u_j
       +x\left(\frac{g}{(1+x)^2}-h\right)v_j\right].
 \label{eq:P11-kH}
\end{equation}
Thus $R_j=0$ is a regular codimension-one level set wherever it
meets $\mathcal U_{\rm eq}$, explicitly the graph $k=k_H$.
A physical point requires $k_H>0$ and
$a_H=(k_H+h)x-(c-g)/(1+x)>0$. Within a strict nonuniform
P05 cell, the portions with $\omega_j\ne0$ are strict candidates,
and this regular surface transfers locally to physical parameter space
along the specified positive branch.

For $u_j>0$, any nonempty intersection $R_j=\omega_j=0$ is
locally a regular codimension-two level set in branch coordinates.
It is codimension one inside $R_j=0$. The same statement transfers
to physical parameters wherever P05's concave-chart hypotheses hold.
For $u_j\le0$ this intersection is empty.
\end{proposition}
\begin{proof}
In \eqref{eq:P11-R} only $-kx$ depends on $k$, proving the
partial derivative and every sign in the solved formula. Strict
selection and $\omega_j\ne0$ are open conditions.
P05 Theorem~\ref{thm:P05-selection} implies $H-G>0$ at any
nonuniform selected representative. On precisely this concave part
of the positive chart, P05 supplies
\[
 \Phi(x,k,c,g,h)=(a,k,c,g,h),\qquad
 \det D\Phi=k+h+\frac{c-g}{(1+x)^2}
     =k+\left(h-\frac{g}{(1+x)^2}\right)
         +\frac{c}{(1+x)^2}>0.
\]
Its local diffeomorphism therefore transfers the regular level set.
No global injectivity or topology after forgetting the branch is claimed.
The derivative $\partial_k R_j$ holds $x,c,g,h$ fixed and lets $a$
follow \eqref{eq:P11-chart}; it is not a derivative along the later
fixed-$a$ continuation path.

For the additional frequency equation, at fixed branch coordinates
\[
 \partial_k\omega_j=0,\qquad
 \partial_g\omega_j=\frac{2xu_j\sin\theta_j}{(1+x)^2}>0
 \quad(u_j>0).
\]
The Jacobian of $(R_j,\omega_j)$ in columns $(k,g)$ has determinant
$-2x^2u_j\sin\theta_j/(1+x)^2\ne0$. The regular level set theorem
proves the local codimension assertion, conditional on existence and
physical positivity; the local diffeomorphism applies under the same
concavity hypotheses just stated. Proposition~\ref{prop:P11-frequency}
excludes all $u_j\le0$ intersections. No unfolding is inferred.
\end{proof}

\begin{proposition}[Exact obstruction and cell realizability criterion]
\label{prop:P11-realizability}
Fix a reachable nonuniform representative $1\le j\le j_{\max}$.
Let $\beta_j=(u_j+u_{j+1})/2$ be P05's midpoint, and put
\begin{equation}
 w_j=1-2u_j^2,\quad L_j^{\rm vert}=\max(0,\beta_j),\qquad
 B_j=\frac{w_j}{4(1-u_j)}.
 \label{eq:P11-realizability-data}
\end{equation}
Its strict P05 cell contains a physical strict spectral Hopf candidate
if and only if
\begin{equation}
 \boxed{w_j>0\quad\hbox{and}\quad L_j^{\rm vert}<B_j.}
 \label{eq:P11-realizability-test}
\end{equation}
More precisely, for a prescribed positive vertex
$z=C/[4(H-G)]$ in that strict cell and prescribed $x>0$, such a
candidate can be realized by some positive physical coefficients if
and only if
\begin{equation}
 \frac{w_j}{4z}>1-u_j+\frac{x}{1+x}.
 \label{eq:P11-fixed-vertex}
\end{equation}
Consequently not every reachable nonuniform cell supports a candidate.
For example $n=8,j=1$ has $w_j=0$ and never reaches $r_j=0$.
For $n=12,j=2$, $u_j=1/2$, $L_j^{\rm vert}=B_j=1/4$;
this cell also has no physical candidate despite $w_j>0$.
\end{proposition}
\begin{proof}
Write $D_c=H-G>0$ and $C=c/(1+x)$ within the selected cell.
At a candidate, \eqref{eq:P11-R} reads
\begin{equation}
 kx=D_cw_j-C\left(1-u_j+\frac{x}{1+x}\right),\qquad
 \frac{kx}{C}=\frac{w_j}{4z}-\left(1-u_j+\frac{x}{1+x}\right).
 \label{eq:P11-k-obstruction}
\end{equation}
Because $k,x,C>0$ and $u_j<1$, this forces $w_j>0$ and
\eqref{eq:P11-fixed-vertex}. It further forces $z<B_j$.
P05's strict cell is
$L_j^{\rm vert}<z<\beta_{j-1}$, where $\beta_{j-1}>0$;
$\beta_j$ exists since $j_{\max}<N$. Thus necessity of
\eqref{eq:P11-realizability-test} follows.

Conversely, if \eqref{eq:P11-realizability-test} holds, choose
\[
 L_j^{\rm vert}<z<\min(\beta_{j-1},B_j).
\]
Both upper bounds exceed the lower bound. Choose $x>0$ sufficiently
small that \eqref{eq:P11-fixed-vertex} holds. For fixed $x,z$
satisfying that inequality, the following construction proves the more
precise sufficiency as well. Choose any $C>0$, set $D_c=C/(4z)$,
and define $k>0$ by \eqref{eq:P11-k-obstruction}. For a free $G>0$,
set
\begin{equation}
 c=C(1+x),\qquad g=\frac{G(1+x)^2}{x},\qquad
 h=\frac{G+D_c}{x},\qquad
 a=kx+D_c-C+G\left(1+\frac{1+x}{x}\right).
 \label{eq:P11-realization}
\end{equation}
These equations give $H-G=D_c$ and exactly
$a=(k+h)x-(c-g)/(1+x)$. Increasing $G$ makes $a>0$
without changing $z$, $k$ or $r_j=0$. The signed frequency bracket
is $2(D_c+2G)u_j-C$. For $u_j\le0$ it is automatically negative;
for $u_j>0$ it vanishes for at most one $G$, which can be avoided
while retaining positivity. Hence a strict candidate is obtained.
The two examples follow by exact substitution of their cosines and
midpoints. They also show why P05's realization of dominance alone
cannot imply realization of a spectral boundary.
\end{proof}

\subsection{An exact five-node physical candidate}
Use the verified P07 constants and spectra, without a new P07 derivation:
\[
 \phi=\frac{1+\sqrt5}{2},\quad
 \sigma=\frac{\sqrt{10-2\sqrt5}}4>0,\quad
 u_1=\frac1{2\phi},\quad u_2=-\frac\phi2.
\]
P07 shows that $\{1,4\}$ is the only potentially dominant nonuniform
pair and that $\omega_1=\sigma[(G+H)-\phi C]$.

\begin{theorem}[Exact five-node spectral Hopf candidate]
\label{thm:P11-five}
For $n=5,q=2$ set
\begin{equation}
 \begin{gathered}
 x_H=1,\quad c=2,\quad g_H=16,\quad h=6,\\
 k_H=\frac{3\phi-4}{2}=\frac{3\sqrt5-5}{4},\qquad
 a_H=\frac{3\phi+22}{2}=\frac{47+3\sqrt5}{4}.
 \end{gathered}
 \label{eq:P11-tuple}
\end{equation}
This is a positive physical tuple with a positive synchronous equilibrium
$x_H\one$. It is a strict spectral candidate in representative $1$.
Its full labeled spectrum is
\begin{align}
 \lambda_0&=-\frac{3\phi+1}{2},\label{eq:P11-five-spectrum}\\
 \lambda_1&=\ii\omega_H,\qquad \lambda_4=-\ii\omega_H,
       &\omega_H&=\sigma(10-\phi)>0,\nonumber\\
 \lambda_2&=-\frac{3\sqrt5}{2}-\ii\sigma(1+10\phi),
       &\lambda_3&=\overline{\lambda_2}.\nonumber
\end{align}
The critical values are each algebraically and geometrically simple.
The ordering is $r_1=0>r_0>r_2$, with exact noncritical gap
\begin{equation}
 \eta_H=\frac{3\phi+1}{2}>0.
 \label{eq:P11-five-gap}
\end{equation}
The equilibrium is the upper of two distinct positive synchronous roots.
The other root is
\begin{equation}
 x_- =\frac{3(2-\phi)}{3\phi+8}\in(0,1).
 \label{eq:P11-lower-root}
\end{equation}
\end{theorem}
\begin{proof}
Since $5/3<\sqrt5<3$, $k_H>0$ and $a_H>0$;
the remaining three coefficients are positive. In the original
canonical field, each component at $\one$ is exactly
\[
 F_i^{\DM}(\one)=a_H-k_H+\frac22-\frac{16}2-6
     =a_H-k_H-13=0.
\]
This verifies the original model residual before using any spectral data.
Direct substitution gives
\begin{equation}
 \begin{gathered}
 s=2,\quad C=1,\quad G=4,\quad H=6,\quad T=G-H=-2,\quad S=G+H=10,\\
 d=-k_H-\frac32=\frac{1-3\phi}{2}.
 \end{gathered}
 \label{eq:P11-five-effective}
\end{equation}
The raw P04 expression independently gives
$d=a_H-2k_H-1/2-8-6=-k_H-3/2$.
P07's exact real-part formulas now yield
\begin{align*}
 r_1&=\frac{1-3\phi}{2}+\frac{1}{2\phi}+\phi=0,\\
 r_0&=\frac{1-3\phi}{2}+1-2=-\frac{3\phi+1}{2},\\
 r_2&=\frac{1-3\phi}{2}-\frac\phi2-\frac1\phi
       =\frac{3-6\phi}{2}=-\frac{3\sqrt5}{2}.
\end{align*}
Here $\phi^{-1}=\phi-1$ and $2\phi-1=\sqrt5$.
Their noncritical difference is
$r_0-r_2=(3\phi-4)/2=k_H>0$. Equivalently,
$-\phi^2<-2<-\phi^{-2}$, so P07's full ordering table gives
the same strict ordering. In particular pair $\{2,3\}$ is excluded.

For the imaginary parts, P07 gives
$\omega_1=\sigma(10-\phi)$ and
$\omega_2=-\sigma(1+10\phi)$. The exact bounds $1<\phi<2$
imply $10-\phi>8>0$, so the critical frequency is nonzero without
decimal evidence. P04 conjugacy fixes the signs at labels $4,3$.
This proves every entry of \eqref{eq:P11-five-spectrum}.
The general theorem proves simplicity and the inverse gates, and the
strict noncritical ordering gives \eqref{eq:P11-five-gap}.

For the branch identification, use only P02's abbreviations and polynomial:
\begin{equation}
 A=k_H+h=\frac{3\phi+8}{2},\quad b=a_H-A=7,\quad
 \mu=a_H+2-16=\frac{3\phi-6}{2}<0.
 \label{eq:P11-branch-data}
\end{equation}
Because $1$ is a root and $A-b=\mu$, that polynomial factors as
\[
 Q(X)=AX^2-bX-\mu=A(X-1)(X-x_-),\qquad x_-=-\mu/A.
\]
The displayed $x_-$ is positive because $\phi<2$, and is less than
one since $3(2-\phi)<3\phi+8$. Thus the discriminant is
$A^2(1-x_-)^2>0$. Together with $b>0$ and $\mu<0$ this checks
P02's two-positive-root regime fully; those two signs alone would
not suffice without the positive discriminant. The larger root is
exactly $x_H=1$. No P10 theorem is used.
\end{proof}

\subsection{Frozen continuation path and exact spectral data}
\label{slot:P11-3}
\begin{proposition}[Local branch, path, vectors and resolvents for P12/P13]
\label{prop:P11-data}
Hold $a=a_H,k=k_H,c=2,h=6$ fixed and freeze the physical path
\begin{equation}
 \boxed{\pi(g)=(a_H,k_H,2,g,6),\qquad g_H=16.}
 \label{eq:P11-path}
\end{equation}
There is a neighborhood of $16$ with a unique smooth positive
synchronous root near $1$, denoted $x_+(g)$, with $x_+(16)=1$.
It continues the upper root. With the project's node ordering $0,\ldots,4$,
\begin{equation}
 \begin{gathered}
 v_H=\frac1{\sqrt5}(1,e^{2\pi\ii/5},e^{4\pi\ii/5},
                         e^{6\pi\ii/5},e^{8\pi\ii/5})^{\mathsf T},\\
 p_H=v_H,\qquad \overline{v_H}=v^{(4)},\qquad
 J_Hv_H=\ii\omega_Hv_H,\quad J_H^*p_H=-\ii\omega_Hp_H,
 \quad p_H^*v_H=1.
 \end{gathered}
 \label{eq:P11-five-vectors}
\end{equation}
The real orthonormal basis $e_c=\sqrt2\Ree v_H$,
$e_s=\sqrt2\Imm v_H$ has block
$\left(\begin{smallmatrix}0&\omega_H\\-\omega_H&0\end{smallmatrix}\right)$.
Both required inverse operators exist and are exactly
\begin{equation}
 \begin{split}
 J_H^{-1}&=\sum_{m=0}^4\lambda_m^{-1}v^{(m)}v^{(m)*},\\
 (2\ii\omega_H I-J_H)^{-1}
   &=\sum_{m=0}^4(2\ii\omega_H-\lambda_m)^{-1}v^{(m)}v^{(m)*}.
 \end{split}
 \label{eq:P11-five-resolvents}
\end{equation}
In particular the mode-0 and mode-2 denominators useful for later work are
\begin{equation}
 \lambda_0=-\eta_H\ne0,\qquad
 2\ii\omega_H-\lambda_2
       =\frac{3\sqrt5}{2}+\ii\sigma(21+8\phi)\ne0.
 \label{eq:P11-denominators}
\end{equation}
\end{proposition}
\begin{proof}
Let $B(X,g)=a_H-(k_H+6)X+(2-g)/(1+X)$.
P02's scalar field is $f(X,g)=XB(X,g)$. P04's uniform identity
$\lambda_0=\partial_X f$ at synchrony, together with $B(1,16)=0$,
gives $B_X(1,16)=\lambda_0=-\eta_H<0$.
The scalar implicit-function theorem supplies the unique smooth root
near $1$. Positivity persists in a sufficiently small neighborhood,
and P02's distinct roots and their strict ordering persist there, so
this is the locally continued upper root, not uniqueness of all positive
roots. The parameter tuple stays positive for $g$ near $16$.
No derivative of $x_+(g)$ is calculated.

Proposition~\ref{prop:P11-vectors} gives all vector identities and
the real block with its stated sign. The spectral theorem and
Theorem~\ref{thm:P11-strict} give the two sums, all of whose denominators
are nonzero. Finally,
$2(10-\phi)+(1+10\phi)=21+8\phi$ verifies the second denominator
in \eqref{eq:P11-denominators}; its real part is positive.
No tensor, nonlinear Fourier selection rule or Taylor-factor convention
is required for these spectral identities.
\end{proof}

\begin{proposition}[Spectral hypothesis ledger at the five-node point]
\label{prop:P11-ledger}
The following is the precise P12/P13 handoff, subject to independent
A11 audit of this DRAFT.
\begin{enumerate}[label=H\arabic*.,leftmargin=3em]
 \item All five physical coefficients in \eqref{eq:P11-tuple} are positive.
 \item $x_H=1$ is an exact positive synchronous equilibrium of the original DM field.
 \item The path varying $g$ in \eqref{eq:P11-path} has the smooth local
       upper branch $x_+(g)$ through $x_+(16)=1$.
 \item Labels $1,4$ give $\pm\ii\omega_H$ with $\omega_H=\sigma(10-\phi)>0$.
 \item Each critical complex eigenvalue is algebraically and geometrically simple.
 \item Every other eigenvalue has real part at most $-\eta_H<0$,
       where $\eta_H=(3\phi+1)/2$.
 \item $0\notin\spec(J_H)$.
 \item $2\ii\omega_H\notin\spec(J_H)$.
 \item Exact normalized right and adjoint vectors are \eqref{eq:P11-five-vectors};
       the pairing and real-plane orientation are fixed.
 \item Spectral transversality is \textbf{not established}; no state or
       eigenvalue derivative along \eqref{eq:P11-path} has been computed.
 \item The first Lyapunov coefficient and nonlinear nondegeneracy are
       \textbf{not established}; no tensor or normal-form coefficient is computed.
 \item No periodic orbit, periodic-branch stability or rotating wave is established.
\end{enumerate}
\end{proposition}
\begin{proof}
Items H1--H2 and H4--H6 are Theorem~\ref{thm:P11-five}; H3 and
H9 are Proposition~\ref{prop:P11-data}; H7--H8 follow from
Theorem~\ref{thm:P11-strict}. H10--H12 explicitly delimit unproved
later conclusions, rather than adding assumptions already certified here.
\end{proof}

\paragraph{Scope, evidence and frozen planning language.}
The exact five-node candidate and the optional cell criterion supply
analytic nonemptiness results, with no numerical premise. Separate
computational checks are corroboration only. The archived decimal
candidate C073 is not used or reproduced. B02 is fixed for this DRAFT
to the upper branch and path \eqref{eq:P11-path}; B06 is met for these
spectral data by exact algebra. B07's pairing and adjoint sign agree
with the current convention, but its remaining Taylor and nonlinear
normalization issues remain for P13. Neither a first loss along the
path nor an open transversality/nonlinear-nondegeneracy regime has been
proved. The accompanying \path{PROPOSED_P11_scope_clarification.md}
records the proposed refinement of older first-crossing wording; all
frozen planning and novelty documents remain unchanged. P12 is not begun.
